% \documentclass[conference]{IEEEtran}
\documentclass[conference,a4paper]{IEEEtran}
\usepackage[a4paper, margin=19mm]{geometry}
\IEEEoverridecommandlockouts

\usepackage{cite}
\usepackage{amsmath,amssymb,amsfonts}
\usepackage{graphicx}
\usepackage{textcomp}
\usepackage{xcolor}
\usepackage{booktabs}
\usepackage{multirow}
\usepackage{microtype}

\begin{document}

\title{MIS: Diagnosing Bottlenecks in Copy-Paste Augmentation
for Agricultural Object Detection}

\author{
    \IEEEauthorblockN{Huy Khanh Tran, Yuichiro Nomura, and Hiroshi Mineno}
    \IEEEauthorblockA{Graduate School of Integrated Science and Technology,
    Shizuoka University, Japan \\
    tran.huy.khanh.25@shizuoka.ac.jp}
}

\maketitle

\begin{abstract}
Copy-paste augmentation achieves strong results for instance
segmentation yet yields inconsistent or negative gains for
bounding-box detection in agricultural scenes.
We propose the \textbf{Minimum Inference Score (MIS)}, a
difficulty metric derived from a trained baseline detector
that correlates strongly with detection failure modes---miss
rate and false-positive rate---across objects and images.
MIS partition analysis reveals whether hard instances
concentrate in constrained regimes---driven by object
size, domain composition, or both.
When such concentration is detected, copy-paste fails
to provide reliable gains across all tested selection
strategies.
Furthermore, by examining which difficulty groups dominate
each MIS partition, practitioners can anticipate which
selection strategies will produce the most harmful
outcomes---all before any augmented training run.
We validate both findings across three agricultural
benchmarks, providing bidirectional validation across
constrained and unconstrained datasets.
\end{abstract}

\begin{IEEEkeywords}
object detection, copy-paste augmentation, difficulty scoring,
agricultural object detection, data augmentation diagnostic
\end{IEEEkeywords}

% ----------------------------------------------------------
\section{Introduction}
Copy-paste augmentation composites cropped instances into
training images to increase diversity~\cite{kisantal2019augmentation,
ghiasi2021simple}, yet yields inconsistent or negative gains
for bounding-box detection in agricultural imagery, forcing
costly multi-seed search without a principled stopping criterion.
This failure persists across all tested selection strategies,
suggesting a dataset-level bottleneck that no sampling policy
can resolve.

We introduce the \textbf{Minimum Inference Score (MIS)}, a
metric computed from a trained baseline detector that captures
model failure patterns on a given dataset.
MIS partition analysis enables two practical outputs without
any augmented run: it determines whether copy-paste can supply
useful training signal, by testing whether hard instances
concentrate in constrained regimes---validated across
MinneApple~\cite{hani2020minneapple},
GWHD~2021~\cite{david2021global}, and
WGISD~\cite{santos2020grape}---and it predicts which selection
strategies risk harmful outcomes, purely from partition
composition.


% ----------------------------------------------------------
\section{Methodology}
\subsection{Scoring Formulation}

Let $G(I)$ denote ground-truth boxes in image $I$.
A trained detector produces predictions $p$ with confidence $c_p$ at
$\tau_{\mathrm{conf}}{=}0.001$ to retain near-miss candidates.
For each $g\in G(I)$, let
$\mathcal{P}(g){=}\{p:\mathrm{IoU}(p,g){\ge}0.5\}$.
The \textit{Object Difficulty Score} is
\begin{equation}
S_{\mathrm{obj}}(g) =
\begin{cases}
  \min_{p \in \mathcal{P}(g)}\!
    [\alpha(1{-}c_p) + \beta(1{-}u_p)],
    & \mathcal{P}(g) \neq \varnothing \\[2pt]
  \alpha + \beta, & \mathcal{P}(g) = \varnothing
\end{cases}
\label{eq:obj}
\end{equation}
with $\alpha{=}\beta{=}0.5$ and $u_p{=}\mathrm{IoU}(p,g)$.
Unmatched objects receive the maximum score $(\alpha{+}\beta{=}1)$.
The \textit{Image Difficulty Score} is
\begin{equation}
S_{\mathrm{img}}(I)
  = \frac{1}{|G(I)|}\sum_{g} S_{\mathrm{obj}}(g)
  + w_2 \cdot \mathrm{FPrate}(I),
\label{eq:img}
\end{equation}
where $\mathrm{FPrate}(I)$ counts unmatched predictions per image,
giving $S_{\mathrm{obj}}\in[0,1]$ per object and $S_{\mathrm{img}}$
per image. By construction,  $S_{\mathrm{img}}$ encodes both
axes of detection failure (miss rate and FP rate); w2 is set
by grid search to balance their contributions.
MIS is model-conditioned, reflecting structural
properties only through detector behavior.

\subsection{Augmentation Protocol}
We train a YOLOv8s baseline and partition all training objects by
$S_{\mathrm{obj}}$ and all training images by $S_{\mathrm{img}}$ into
equal tertiles (Low, Medium, High).
Four copy-paste conditions are evaluated:
\textbf{Score High} pastes objects from the highest-$S_{\mathrm{obj}}$
tertile into the highest-$S_{\mathrm{img}}$ images;
\textbf{Score Medium} and \textbf{Score Low} follow the same logic
for their respective tertiles.
\textbf{Random} selects both destination images and source objects
without score-based filtering.
This design isolates whether targeting specific score regimes
affects outcome. Training data is expanded by 30\% via unblended within image
copy-paste~\cite{kisantal2019augmentation}: 2--8 objects per image,
scale~$[0.9,1.1]$, rotation~$\pm15^\circ$, $3{\times}$~maximum
reuse; all other settings follow default YOLOv8 protocol.

% ----------------------------------------------------------
\section{Results}
\subsection{MIS Reflects Detection Failure}

Pearson correlations between $S_{\mathrm{img}}$ and miss rate
on the validation set are $r{=}0.79$, $0.81$, $0.91$ on
MinneApple, GWHD, and WGISD; correlations with false positive
rate are $r{=}0.79$, $0.81$, $0.84$. These confirm that images
with high $S_{\mathrm{img}}$ are genuinely harder to
detect---the model misses more objects and produce more false positives on 
them---validating
MIS as a reliable difficulty proxy.

\subsection{MIS Identifies Dataset Failure Sources}

\textbf{Size (MinneApple).}
$S_{\mathrm{obj}}$ strongly correlates with object size
(Spearman~$r_s{=}{-}0.82$): 64.4\% of the top-difficulty
tertile falls in \textit{Very Small} instances
(${\le}400$\,px$^2$, low AP), indicating the model struggles
primarily with small objects.

\textbf{Domain (GWHD~2021).}
Per-domain median MIS strongly correlates with
domain AP (Pearson $r{=}{-}0.90$; Fig.~\ref{fig:gwhd}): MIS
assigns higher scores to harder acquisition domains, purely
from detector behavior and without access to domain labels.

\begin{table}[t]
\centering
\caption{MinneApple: $S_{\mathrm{obj}}$ captures size-level difficulty.}
\label{tab:size_val}
\footnotesize
\setlength{\tabcolsep}{3.5pt}
\renewcommand{\arraystretch}{0.95}
\begin{tabular}{lccc}
\toprule
\textbf{Category} & \textbf{AP} & \textbf{Mean $S_{\mathrm{obj}}$}
  & \textbf{Top tertile} \\
\midrule
Very Small ($\le$400\,px$^2$) & 0.071 & 0.152 & 64.4\% \\
Small (400--1024\,px$^2$)     & 0.270 & 0.076 & 33.9\% \\
Medium ($>$1024\,px$^2$)      & 0.493 & 0.053 &  1.7\% \\
\bottomrule
\end{tabular}
\end{table}

\begin{figure}[t]
  \centering
  \includegraphics[width=\columnwidth]{gwhd_domain_score.png}
  \caption{GWHD~2021: per-domain AP vs.\ median image difficulty score.}
  \label{fig:gwhd}
\end{figure}

\subsection{Application: Copy-Paste Diagnostic and Strategy Risk}

\textbf{MinneApple (resolution ceiling).}
We define \emph{resolution ceiling} as a size below which an object's
footprint on the detector's feature maps is too small to
localize reliably, no matter how many training examples exist.
We call this a \emph{structural constraint}: it is a limit of
representation capacity, not of data quantity.
MIS partition analysis shows 64.4\% of high-difficulty objects
fall in the Very Small category (${\le}400$\,px$^2$, AP\,${=}$\,0.071).
This signals a resolution ceiling: pasting additional 
objects provides no useful learning signal because the model
already fails on them structurally, not from lack of examples.
This is confirmed: all copy-paste conditions yield negative
$\Delta$AP (Table~\ref{tab:outcomes}), regardless of selection
strategy, with no condition reversing the outcome direction.


\textbf{GWHD~2021 (domain difficulty and strategy risk).}
MIS partitioning reveals near-complete domain separation:
low-difficulty images are 66\%  ETHZ\_1/Rres\_1
(AP\,${\approx}$\,0.65), while high-difficulty images are
78\% Arvalis variants (AP\,${\approx}$\,0.46).
Since difficulty stems from each domain's image-level
properties rather than a lack of instances, copy-paste cannot
help: it adds instances but cannot change how a domain looks
to the detector. This is confirmed: all conditions yield
negative $\Delta$AP.
Partition composition further predicts per-strategy risk from
the domains each strategy sources from: Score~Low draws
exclusively from high-AP domains, flooding training with
redundant patterns (consistent drop, $-0.034$,
$\sigma{=}0.005$); Score~High draws from hard Arvalis domains,
introducing high-variance signal the model cannot reliably
exploit (seed instability, $-0.006$, $\sigma{=}0.040$); Random
and Medium source from mixed partitions, giving moderate,
more stable drops ($-0.015$/$-0.017$, $\sigma{=}0.023$/$0.021$).

\textbf{WGISD (positive control).}
MIS raises no flag: difficulty is distributed across size
categories without concentration, and no domain separation is
present, signaling that copy-paste can introduce useful
learning signal. This is confirmed: all conditions yield
positive $\Delta$AP (Table~\ref{tab:outcomes}), with
Score~Medium achieving the best result.

\begin{table}[t]
\centering
\caption{Copy-paste outcomes (3-seed mean $\pm\,\sigma$).
Bottleneck:  failure source identified by MIS.}
\label{tab:outcomes}
\footnotesize
\setlength{\tabcolsep}{2.5pt}
\renewcommand{\arraystretch}{0.95}
\begin{tabular}{lllll}
\toprule
\textbf{Dataset}
  & \textbf{Bottleneck}
  & \textbf{Base AP}
  & \textbf{Best $\Delta$}
  & \textbf{Worst $\Delta$} \\
\midrule
MinneApple
  & Size
  & $0.353\pm0.005$
  & $-0.003$ (Hi)
  & $-0.008$ (Rnd) \\
GWHD~2021
  & Domain
  & $0.229\pm0.013$
  & $-0.006^{*}$ (Hi)
  & $-0.034$ (Lo) \\
WGISD
  & ---
  & $0.480\pm0.007$
  & $+0.009$ (Med)
  & $+0.001$ (Hi) \\
\bottomrule
\multicolumn{5}{l}{\scriptsize
  $^*$Score High on GWHD: $\sigma{=}0.040$, high seed variance.}
\end{tabular}
\end{table}

\subsection{Limitations}
Two aspects remain unvalidated: $w_2$ in Eq.~\eqref{eq:img}
was fixed by grid search on one configuration, so its effect on
the tertile partitions and diagnostic conclusions is untested;
and all results use a single detector, YOLOv8s, so transfer to
other architectures is untested. Both are left to future work,
via sensitivity analysis and cross-architecture validation.

% ----------------------------------------------------------
\section{Conclusion}

Copy-paste augmentation fails when hard instances concentrate
in structurally constrained regimes---below the detector's
resolution range or within acquisition-defined domains that
instance-level augmentation cannot bridge.
MIS makes these constraints visible before any augmented run:
on MinneApple and GWHD~2021, partition analysis correctly
predicted failure across all four selection strategies;
on WGISD, the absence of such concentration correctly
predicted positive gains across all conditions---providing
bidirectional validation that MIS functions as a diagnostic,
not merely a failure detector.
On domain-structured datasets, MIS partition composition
delivers a second output: it identifies risky selection
strategies from partition composition alone, before any
augmented run. Across our experiments, 24 augmented runs on MIS-flagged
datasets yielded zero reliable gains, a compute cost MIS
avoids entirely. These findings hold for the single weighting
and detector validated here (Section~III-D); future work
includes guiding instance-level selection when constraints are
absent, and pairing MIS with generative methods for hard
regimes copy-paste cannot reach.

\section*{Acknowledgment}
This work was supported by JST FOREST JPMJFR201B,
JSPS KAKENHI JP25H01110.

\begin{thebibliography}{00}

\bibitem{kisantal2019augmentation}
M.~Kisantal \textit{et~al.},
``Augmentation for small object detection,''
\textit{arXiv:1902.07296}, 2019.

\bibitem{ghiasi2021simple}
G.~Ghiasi \textit{et~al.},
``Simple copy-paste is a strong data augmentation method for
instance segmentation,''
in \textit{Proc.\ CVPR}, 2021, pp.~2918--2928.

\bibitem{hani2020minneapple}
N.~H\"{a}ni, P.~Roy, and V.~Isler,
``MinneApple: A benchmark dataset for apple detection and
segmentation,''
\textit{IEEE RA-L}, vol.~5, no.~2, pp.~852--858, 2020.

\bibitem{david2021global}
E.~David \textit{et~al.},
``Global wheat head detection 2021,''
\textit{Plant Phenomics}, 2021.

\bibitem{santos2020grape}
T.~T.~Santos \textit{et~al.},
``Grape detection, segmentation, and tracking using deep
neural networks,''
\textit{Comput.\ Electron.\ Agric.}, vol.~170, p.~105247, 2020.

\end{thebibliography}

\end{document}
